% !PRINTED 296
\weberpage{296}
%<<CONT
auf dem Gedanken, der dem Charakteristikenbegriff zu Grunde
liegt\,\footnote{\gesperrt{Gauss}, Demonstratio nova theorematis omnem functionem algebraicam rationalem integram unius variabilis in factores reales primi vel
secundi gradus resolvi posse (1799). Beiträge zur Theorie der algebraischen
Gleichungen (1849). \gesperrt{Kronecker}, Monatsberichte der Berliner Akademie,
21. Februar 1878.}. Wir schliessen daher diesen Abschnitt passend mit einer
Darlegung dieses Beweises.

Es sei also $z = x + yi$ und
\weq{1}{F(z) = \varphi(x, y) + i\,\psi(x, y)}
eine ganze rationale Function von $z$ vom $n^{\text{ten}}$ Grade mit reellen
oder complexen Coëfficienten, in der $z^n$ den Coëfficienten 1 hat,
also
\weq{2}{F(z) = z^n + a_1 z^{n-1} + a_2 z^{n-2} + \cdots + a_n.}

Wir setzen wieder voraus, dass $F(z)$ und $F'(z)$ keinen gemeinsamen Theiler haben, und der Beweis, dass $F(z)$ für $n$
Werthe von $z$ verschwindet, beruht nun darauf, dass, wenn wir
für $f$ einen Kreis von hinlänglich grossem Radius wählen, die
Charakteristik des Functionensystems $f, \varphi, \psi$ direct bestimmt
werden kann und sich gleich $n$ findet. Dann folgt nach \S.\,96,
dass in dem so gewählten Kreise $n$ Punkte $z$ liegen, in denen
$F(z)$ verschwindet.

Wir bedienen uns der Polarcoordinaten, und setzen, da
$a_1, a_2, a_3 \dots a_n$ auch complex sein können
\[
a_1 = p_1 e^{iq_1},\quad a_2 = p_2 e^{iq_2}, \dots a_n = p_n e^{iq_n}
\]
\[
z = R e^{i\vartheta},
\]
worin $p_1, p_2, \dots p_n, R$ positiv sind und die Winkel $q_1, q_2, \dots q_n, \vartheta$
in irgend einem Intervall vom Umfang $2\pi$ gewählt werden können.
Hierdurch wird
\weq{3}{
\begin{aligned}
\varphi(x, y) &= R^n \cos n\vartheta + p_1 R^{n-1} \cos[(n-1)\vartheta + q_1] \\
&\quad + \cdots + p_n \cos q_n \\[4pt]
\psi(x, y) &= R^n \sin n\vartheta + p_1 R^{n-1} \sin[(n-1)\vartheta + q_1] \\
&\quad + \cdots + p_n \sin q_n,
\end{aligned}
}
und wir bilden auch noch die nach $\vartheta$ genommenen Ableitungen
%CONT>>
